\subsection{Уравнение кинематической связи (ниточка, палочка)}

%%%%%%%%%%%%%%%%%%%%%%%%%%%%%%%%%%%%%%%%%%%%%%%%%%%%
%%%%%%%%%%%%%%%%%%%%%%%%%%%%%%%%%%%%%%%%%%%%%%%%%%%%
 
 \z[{\begin{wrapfigure}[6]{r}{0.3\textwidth}%
		\vspace{-10pt}%
		\begin{flushright}%
 			\includegraphics[width=0.25\textwidth]{1.5.1}%
		\end{flushright}%
	\end{wrapfigure}%
	}]%
	{На наклонную плоскость, составляющую с горизонтом угол $\alpha$, опирается стержень, который может перемещаться только по вертикали благодаря направляющему устройству $AB$ (см. рисунок). С какой скоростью v поднимается стержень, если наклонная плоскость влево в горизонтальном направлении со скоростью $u$?}
%
{$v=\tan{\alpha}$}

%%%%%%%%%%%%%%%%%%%%%%%%%%%%%%%%%%%%%%%%%%%%%%%%%%%%

 \z[{\begin{wrapfigure}[6]{r}{0.3\textwidth}%
		\vspace{-25pt}%
		\begin{flushright}%
 			\includegraphics[width=0.3\textwidth]{1.5.2}%
		\end{flushright}%
	\end{wrapfigure}%
	}]%
	{Стержень длиной $L$ шарнирно соединен с муфтами $A$ и $B$, которые перемещаются по двум взаимно перпендикулярным рейкам (см. рисунок). Муфта $А$ движется вдоль оси $Ox$ с постоянной скоростью $v$. Приняв за начало отсчета времени момент, когда муфта $A$ находилась в точке $O$, определите зависимость от времени координаты $y(t)$ и скорости $v_y (t)$ муфты $B$.}
%
{$y(t)=\sqrt{L^2-v^2t^2}$;~$v_y(t)=\frac{v^2t^2}{\sqrt{L^2-v^2t^2}}$}
	
%%%%%%%%%%%%%%%%%%%%%%%%%%%%%%%%%%%%%%%%%%%%%%%%%%%%
\z[{\begin{wrapfigure}[6]{r}{0.3\textwidth}%
		\vspace{-10pt}%
		\begin{flushright}%
 			\includegraphics[width=0.3\textwidth]{1.5.3}%
		\end{flushright}%
	\end{wrapfigure}%
	}]%
	{Небольшой брусок через систему блоков связан нерастяжимой нитью с длинной тележкой, которая может катиться по горизонтальной поверхности. Брусок кладут на тележку и приводят в движение с постоянной скоростью $v = 2\,\nicefrac{\text{м}}{\text{с}}$, направленной горизонтально вдоль тележки (см. рис.). Какую скорость относительно бруска будет иметь тележка в тот момент, когда угол между наклонной нитью и горизонтом составит $\alpha = 60\degree$? Считайте, что в указанный момент тележка не доехала до стены, к которой прикреплены блоки.}
%
{$v_\text{отн}=v(1+\cos{\alpha})=3\,\nicefrac{\text{м}}{\text{с}}$}

%%%%%%%%%%%%%%%%%%%%%%%%%%%%%%%%%%%%%%%%%%%%%%%%%%%%	
\z[{\begin{wrapfigure}[4]{r}{0.3\textwidth}%
		\vspace{-20pt}%
		\begin{flushright}%
 			\includegraphics[width=0.3\textwidth]{1.5.4}%
		\end{flushright}%
	\end{wrapfigure}%
	}]%
	{Двое рабочих передвигают тяжёлый ящик с помощью канатов (см. рисунок). Найти скорость ящика в момент, когда канаты образуют угол $\alpha$, а скорости рабочих $v_1$ и $v_2$ направлены вдоль канатов.}
%
{$v=\frac{1}{\sin \alpha}\sqrt{v_1^2+v_2^2-2v_1v_2\cos \alpha}$}

%%%%%%%%%%%%%%%%%%%%%%%%%%%%%%%%%%%%%%%%%%%%%%%%%%%%
 \z{По двум пересекающимся под углом $\alpha = 30\degree$ дорогам движутся к перекрёстку два автомобиля: один со скоростью $v_1  = 10\,\nicefrac{\text{м}}{\text{с}}$, второй --- с $v_2=10\sqrt{3}\approx17.3\,\nicefrac{\text{м}}{\text{с}}$. Когда расстояние между автомобилями было минимальным, первый из них находился на расстоянии $S_1  = 200\, \text{м}$ от перекрёстка. На каком расстоянии $S_2$ от перекрёстка в это время находился второй автомобиль?}
%
{$S_2\approx 155\, \text{м}$}

